\chapter{Uzupełnienia węzłów, geometria i chirurgia Dehna}
\label{ch:kt-chirurgia-3d}

W tym rozdziale przechodzimy od krzywych zanurzonych w $S^3$ do
trójrozmaitości z brzegiem. Kluczowym obiektem będzie uzupełnienie
węzła, a kluczową operacją chirurgia Dehna, która pozwala budować
zamknięte trójrozmaitości przez ponowne doklejenie torusów pełnych.
Dzięki temu teoria węzłów staje się narzędziem konstrukcyjnym
w topologii trójwymiarowej.

\section{Otoczenie tubularne, meridian i długość}
\label{sec:kt-peryferia}

\begin{definicja}[Zewnętrze węzła]
Niech $K\subset S^3$ będzie węzłem. Ustalmy domknięte otoczenie
tubularne $N(K)\cong S^1\times D^2$, którego rdzeniem jest $K$.
\emph{Zewnętrzem węzła} nazywamy zwartą trójrozmaitość
$X_K:=S^3\setminus\operatorname{int}N(K)$. Jej brzeg
$\partial X_K=\partial N(K)\cong T^2$ jest torusem.
\end{definicja}

\begin{definicja}[Meridian i długość preferowana]
Niech $K\subset S^3$ będzie zorientowanym węzłem.
\emph{Meridianem} $\mu\subset\partial X_K$ nazywamy
zorientowaną krzywą prostą, która ogranicza dysk
w $N(K)$ i spełnia $\operatorname{lk}(\mu,K)=1$.
\emph{Długością preferowaną} (\emph{preferred longitude})
$\lambda\subset\partial X_K$ nazywamy zorientowaną,
prostą krzywą o liczbie splecenia zero z $K$, reprezentującą
w $H_1(N(K);\mathbb Z)$ tę samą klasę co zorientowany rdzeń.
Klasę $[\lambda]$ można równoważnie określić jako generator
jądra $H_1(\partial X_K;\mathbb Z)\to H_1(X_K;\mathbb Z)$,
z orientacją zgodną z $K$.
\end{definicja}

\begin{stwierdzenie}[Homologia peryferyjna]
\label{prop:kt-peripheral-homology}
Dla każdego zorientowanego węzła $K\subset S^3$ klasy
$[\mu],[\lambda]$ tworzą bazę $H_1(\partial X_K;\mathbb Z)
\cong\mathbb Z^2$, przy czym odwzorowanie
$i_*\colon H_1(\partial X_K;\mathbb Z)\to H_1(X_K;\mathbb Z)$
spełnia $i_*[\mu]=1$ i $i_*[\lambda]=0$, po identyfikacji
$H_1(X_K;\mathbb Z)\cong\mathbb Z$ przez liczbę splecenia.
\end{stwierdzenie}

\begin{proof}
W identyfikacji $N(K)\cong S^1\times D^2$ meridian
jest brzegiem dysku $\{p\}\times D^2$, nie krzywą
przecinającą ten dysk poprzecznie. Wybierzmy podłużną
krzywą $\lambda_0=S^1\times\{q\}$, z $q\in\partial D^2$.
Klasy $[\mu],[\lambda_0]$ tworzą bazę homologii torusa
jako klasy dwóch czynników iloczynu $S^1\times S^1$.
Oznaczmy przez $j$ włączenie brzegu do $N(K)$.

Stosujemy ciąg Mayera--Vietorisa do rozkładu
$S^3=X_K\cup N(K)$, po pogrubieniu obu części
kołnierzem brzegu, aby otrzymać otwarte pokrycie
o przecięciu równoważnym homotopijnie torusowi.
Ponieważ $H_2(S^3;\mathbb Z)=H_1(S^3;\mathbb Z)=0$,
otrzymujemy izomorfizm
\[
(i_*,-j_*)\colon H_1(\partial X_K;\mathbb Z)
\xrightarrow{\ \cong\ }
H_1(X_K;\mathbb Z)\oplus H_1(N(K);\mathbb Z).
\]
Jądro $j_*$ jest generowane przez $[\mu]$, gdyż
$j_*[\lambda_0]=[K]$ i $j_*[\mu]=0$.
Ograniczenie powyższego izomorfizmu do tego jądra
daje izomorfizm $\mathbb Z[\mu]\cong H_1(X_K;\mathbb Z)$.
Liczba splecenia z $K$ jest homomorfizmem na tej grupie
i przyjmuje wartość $1$ na meridianie, zatem jest
wskazaną identyfikacją z $\mathbb Z$.

Każda klasa przechodząca przez $j_*$ na $[K]$ ma
postać $[\lambda_0]+k[\mu]$. Jej liczba splecenia
wynosi $\operatorname{lk}(\lambda_0,K)+k$.
Wybór $k=-\operatorname{lk}(\lambda_0,K)$ daje
klasę $[\lambda]$ o liczbie splecenia zero.
Współrzędne $(k,1)$ są względnie pierwsze, więc
klasa ma prostego reprezentanta na torusie:
w modelu $\mathbb R^2/\mathbb Z^2$ jest nim
obraz prostej $t\mapsto (kt,t)$, $t\in\mathbb R/\mathbb Z$.
Zmiana bazy z $([\mu],[\lambda_0])$ na
$([\mu],[\lambda])$ ma wyznacznik $1$.
W szczególności $i_*[\mu]=1$, $i_*[\lambda]=0$.
\end{proof}

\begin{stwierdzenie}[Jednoznaczność długości preferowanej]
Niech $K\subset S^3$ będzie zorientowanym węzłem, $N(K)$ jego
otoczeniem tubularnym, a $X_K=S^3\setminus\operatorname{int}N(K)$.
Istnieje dokładnie jedna klasa $\lambda\in H_1(\partial X_K;\mathbb Z)$,
która przechodzi na generator $[K]\in H_1(N(K);\mathbb Z)$
i jednocześnie znika w $H_1(X_K;\mathbb Z)$.
W szczególności długość preferowana nie zależy od ramowania
wybranego przy początkowej identyfikacji otoczenia tubularnego.
\end{stwierdzenie}

\begin{proof}
Wybierzmy dowolną krzywą podłużną $\lambda_0$ przechodzącą
na $[K]$ w $H_1(N(K))$. Każda inna klasa o tej własności
ma postać $\lambda_0+k\mu$ dla pewnej liczby całkowitej $k$,
bo jądro homomorfizmu $H_1(\partial N(K))\to H_1(N(K))$
jest generowane przez meridian. Odwzorowanie
$H_1(X_K)\to\mathbb Z$, $[\gamma]\mapsto
\operatorname{lk}(\gamma,K)$, jest izomorfizmem.
Istnieje zatem dokładnie jedna wartość $k$, dla której
$\operatorname{lk}(\lambda_0+k\mu,K)=0$: mianowicie
$k=-\operatorname{lk}(\lambda_0,K)$.
Ponieważ $\operatorname{lk}(\mu,K)=1$, otrzymana
klasa jest jedyną długością preferowaną.
\end{proof}

\begin{center}
\begin{tikzpicture}[scale=.91,line cap=round,>=Stealth]
  \fill[ktBlue!14] (0,0) ellipse (3 and 1.32);
  \draw[ktBlue!70!black,line width=1.25pt] (0,0) ellipse (3 and 1.32);
  \fill[white] (0,0) ellipse (1.15 and .48);
  \draw[ktInk,line width=1.15pt] (0,0) ellipse (1.15 and .48);
  \draw[ktTeal,line width=3.1pt,densely dashed,opacity=.7] (-2.09,0)
        arc[start angle=180,end angle=360,x radius=2.09,y radius=.76];
  \draw[ktTeal,line width=3.2pt,-{Stealth[length=3mm]}] (2.09,0)
        arc[start angle=0,end angle=180,x radius=2.09,y radius=.76];
  \draw[ktCoral,line width=3.1pt,-{Stealth[length=2.5mm]}]
        (2.08,0) ellipse (.43 and .62);
  \node[font=\small\bfseries,text=ktTeal] at (-.15,1.75) {$\lambda$: kierunek podłużny};
  \draw[ktTeal,->,thick] (1.15,1.5)--(1.12,.64);
  \node[font=\small\bfseries,text=ktCoral] at (3.68,.68) {$\mu$};
  \draw[ktCoral,->,thick] (3.42,.55)--(2.45,.37);
  \node[text=ktInk,font=\footnotesize] at (0,-1.75)
        {Schemat torusa $\partial N(K)$: dwa kierunki homologii $H_1(T^2;\mathbb Z)$.};
\end{tikzpicture}
\end{center}

\begin{uwaga}
Na rysunku długość i meridian pokazano schematycznie;
przerywana część długości oznacza fragment po drugiej
stronie powierzchni. Wybieramy identyfikację produktu
$N(K)\cong S^1\times D^2$ dostosowaną do długości
preferowanej. W modelu ilorazowym torusa ich dokładnymi
reprezentantami są krzywe $(t,0)$ i $(0,t)$ w
$\mathbb R^2/\mathbb Z^2$ (pierwsza współrzędna
odpowiada rdzeniowi). Przecinają się one dokładnie raz
na torusie; dodatkowe przecięcia na schemacie mogą
wynikać z rzutu na płaszczyznę.
\end{uwaga}

\section{Nieściśliwe torusy i topologia uzupełnienia}
\label{sec:kt-torus-jsj}

\begin{definicja}[Powierzchnia nieściśliwa]
Niech $M$ będzie zwartą, orientowalną trójrozmaitością,
a $F\subset M$ spójną, dwustronną, właściwie zanurzoną
powierzchnią (lub składową brzegu $M$), niebędącą sferą.
\emph{Dyskiem kompresującym} nazywamy zanurzony dysk
$D\subset M$ taki, że $D\cap F=\partial D$, wnętrze
$D$ leży we wnętrzu $M$, a $\partial D$ jest prostą
krzywą zamkniętą leżącą we wnętrzu powierzchni $F$ i reprezentującą
nietrywialną klasę w $\pi_1(F)$. Powierzchnia $F$ jest
\emph{nieściśliwa}, jeżeli nie ma dysku kompresującego.
Dla powierzchni niespójnej stosujemy ten warunek
składowa po składowej.
\end{definicja}

\begin{uwaga}
Inkluzja $\pi_1(F)\to\pi_1(M)$ jest injektywna wtedy
i tylko wtedy, gdy $F$ jest nieściśliwa. Kierunek
od injektywności wynika natychmiast z definicji:
brzeg dysku kompresującego jest nietrywialny w
$\pi_1(F)$, ale trywialny w $\pi_1(M)$. Kierunek
odwrotny jest konsekwencją twierdzenia o pętli
Papakyriakopoulosa (dla powierzchni wewnętrznej
stosujemy je po rozcięciu wzdłuż $F$, wraz z twierdzeniem
van Kampena). Przyjmujemy ten klasyczny wynik z
topologii trójrozmaitości; zob. C. D. Papakyriakopoulos,
\emph{On Dehn's lemma and the asphericity of knots},
Ann. of Math. 66 (1957), 1--26.
\end{uwaga}

\begin{stwierdzenie}[Nieściśliwość brzegu zewnętrza]
\label{prop:kt-boundary-incompressible}
Jeśli $K\subset S^3$ jest nietrywialnym węzłem, to
$\partial X_K$ jest nieściśliwy w $X_K$.
Dla węzła trywialnego jego zewnętrze jest torusem pełnym,
a torus brzegowy jest ściśliwy.
\end{stwierdzenie}

\begin{proof}
Przypuśćmy, że torus $\partial X_K$ jest ściśliwy i niech
$D\subset X_K$ będzie dyskiem kompresującym, to znaczy
$\partial D$ jest istotną krzywą na $\partial X_K$.
Każda istotna prosta krzywa zamknięta na torusie
reprezentuje klasę pierwotną; korzystamy z klasyfikacji
krzywych prostych na torusie. Zatem jej klasa ma
postać $p[\mu]+q[\lambda]$ z $\gcd(p,q)=1$.
Ponieważ $[\partial D]=0$ w $H_1(X_K)$,
ze stwierdzenia \ref{prop:kt-peripheral-homology}
wynika $p=0$, więc $q=\pm1$.
Po zmianie orientacji brzegu dysku możemy przyjąć
$q=1$. Krzywa $\partial D$ jest izotopijna na
torusie do reprezentanta klasy $[\lambda]$.

W początkowych współrzędnych produktu
$N(K)\cong S^1\times D^2$ zapiszmy
$[\lambda]=[\lambda_0]+k[\mu]$.
Dyfeomorfizm
\[
\Phi_k(e^{\mathrm it},z)=(e^{\mathrm it},e^{\mathrm ikt}z)
\]
ustala rdzeń i przekształca podłużną krzywą
$S^1\times\{1\}$ w krzywą klasy $[\lambda]$.
W szczególności nie twierdzimy, że różne nachylenia
są izotopijne na torusie brzegowym. Obraz pierścienia
$S^1\times[0,1]\subset S^1\times D^2$ przez
$\Phi_k$ ma parametryzację
$(e^{\mathrm it},r e^{\mathrm ikt})$, $0\leq r\leq1$,
i łączy rdzeń z reprezentantem długości preferowanej.
Izotopię tego reprezentanta do $\partial D$ na
brzegu przedłużamy na kołnierz brzegu $N(K)$,
pozostawiając rdzeń nieruchomy. Otrzymujemy
zanurzony pierścień $A$ o brzegu $K\cup\partial D$
i wnętrzu w $\operatorname{int}N(K)\setminus K$.

Wnętrza $A$ i $D$ są rozłączne, gdyż leżą po
przeciwnych stronach torusa. Sklejenie $A\cup D$
wzdłuż wspólnej składowej brzegu jest dyskiem:
doklejenie pierścienia do dysku jedynie przedłuża
jego kołnierz. Po wygładzeniu narożnika otrzymujemy
gładko zanurzony dysk w $S^3$ o brzegu $K$.
Korzystamy z klasycznego kryterium węzła trywialnego:
węzeł ograniczający zanurzony dysk w $S^3$ jest
trywialny (regularne otoczenie dysku jest kulą,
w której jego brzeg można sprowadzić do standardowego
okręgu). Otrzymujemy sprzeczność.

Odwrotnie, dla węzła trywialnego jego zewnętrze jest
torusem pełnym. Dysk południkowy tego torusa stanowi
dysk kompresujący $\partial X_U$.
\end{proof}

\begin{definicja}[Węzeł satelitarny]
Węzeł $K\subset S^3$ nazywamy \emph{satelitarnym}, jeżeli jego
zewnętrze $X_K$ zawiera zanurzony torus
$T\subset\operatorname{int}X_K$, który jest nieściśliwy
i nie jest równoległy do brzegu, to znaczy nie ogranicza
z $\partial X_K$ obszaru homeomorficznego z
$T^2\times[0,1]$ z tymi torusami jako składowymi brzegu.
Taki torus nazywamy istotnym. Geometryczna konstrukcja
satelity polega na umieszczeniu wzoru w torusie pełnym
o nietrywialnym rdzeniu. Sama taka inkluzja nie wystarcza:
wzór nie może leżeć w kuli zawartej w torusie pełnym
ani być izotopijny do jego rdzenia.
\end{definicja}

\begin{uwaga}
Istnienie torusa nieściśliwego w uzupełnieniu nie musi
wynikać z tego, że sam węzeł jest torusowy. Dla węzłów
torusowych $T(p,q)$, $p,q\geq2$, $\gcd(p,q)=1$,
zewnętrze jest przestrzenią Seiferta nad dyskiem
z dwoma punktami wyjątkowymi i jest atoroidalne.
Korzystamy tu ze standardowej klasyfikacji
nieściśliwych torusów w przestrzeniach Seiferta,
a nie z samego istnienia fibracji Seiferta.
Zob. J. S. Purcell, \emph{Hyperbolic Knot Theory},
AMS, 2020, część poświęcona powierzchniom istotnym
i klasyfikacji zewnętrzy węzłów. Torusy otaczające
węzły towarzyszące odgrywają rolę dla węzłów satelitarnych.
\end{uwaga}

\section{Twierdzenie Gordona--Lueckego}
\label{sec:kt-gl}

\begin{twierdzenie}[Gordon--Luecke; węzły są wyznaczone przez uzupełnienia]
Niech $K_1,K_2\subset S^3$ będą węzłami. Jeżeli istnieje
homeomorfizm zachowujący orientację
$h\colon X_{K_1}\to X_{K_2}$, to istnieje dyfeomorfizm
$F\colon S^3\to S^3$ zachowujący orientację i przenoszący
$K_1$ na $K_2$. W konsekwencji węzły są równoważne
jako węzły niezorientowane. Stwierdzenie nie rozstrzyga
samo przez się równoważności ustalonych orientacji $K_i$.
\end{twierdzenie}

\begin{uwaga}
Przyjmujemy jako wynik zewnętrzny twierdzenie
C. McA. Gordona i J. Lueckego,
\emph{Knots are determined by their complements},
J. Amer. Math. Soc. 2 (1989), 371--415.
Jego postać zachowująca orientację daje dokładnie
powyższe stwierdzenie; w kategorii trójrozmaitości
korzystamy dodatkowo z twierdzenia Moise'a o zgodności
kategorii topologicznej, PL i gładkiej, wraz z
wygładzeniem homeomorfizmu par w przypadku gładkich
węzłów. Głównym składnikiem dowodu Gordona--Lueckego
jest twierdzenie, że wypełnienie Dehna zewnętrza
nietrywialnego węzła dające $S^3$ ma nachylenie
meridionalne. Istotnie, homeomorfizm zewnętrzy
przenosi nachylenie meridionalne na nachylenie
dające $S^3$ na drugim zewnętrzu, a więc zachowuje
nachylenie meridionalne. Taki homeomorfizm brzegu
przedłuża się przez torusy pełne z zachowaniem
rdzeni. Przypadek węzła trywialnego rozpoznajemy
po ściśliwym brzegu na mocy stwierdzenia
\ref{prop:kt-boundary-incompressible}.
Nie wystarcza sam fakt, że $X_K$ ma brzeg torusowy:
przy dowolnej trójrozmaitości z takim brzegiem różne
wypełnienia mogą prowadzić do tej samej przestrzeni.
Specjalna rola $S^3$ jest tu zasadnicza. Twierdzenie
pokazuje, że problem klasyfikacji węzłów można przenieść
do topologii ich zewnętrzy.
\end{uwaga}

\section{Geometria hiperboliczna uzupełnień}
\label{sec:kt-hyperbolic}

\begin{definicja}[Węzeł hiperboliczny]
Węzeł $K\subset S^3$ jest \emph{hiperboliczny}, jeżeli
$S^3\setminus K$ dopuszcza zupełną metrykę riemannowską
o stałej krzywiźnie sekcyjnej $-1$ i skończonej objętości.
Objętość tę oznaczamy $\operatorname{Vol}(K)$.
\end{definicja}

\begin{twierdzenie}[Hiperbolizacja uzupełnień węzłów; Thurston]
Każdy nietrywialny węzeł $K\subset S^3$ należy do jednej
z trzech rozłącznych klas: torusowych, satelitarnych albo
hiperbolicznych. W ostatnim przypadku uzupełnienie posiada
zupełną, skończenie objętościową strukturę hiperboliczną.
\end{twierdzenie}

\begin{uwaga}
Przyjmujemy tu zewnętrzne twierdzenie Thurstona
o hiperbolizacji zewnętrzy węzłów, w postaci trójpodziału
dla nietrywialnych węzłów w $S^3$; zob. J. S. Purcell,
\emph{Hyperbolic Knot Theory}, AMS, 2020, część
poświęcona powierzchniom istotnym i hiperbolizacji.
Nie ograniczamy się do węzłów pierwszych: pojęcie
satelitarności przyjęte powyżej obejmuje również
nietrywialne sumy spójne węzłów.
Dla węzła trywialnego zewnętrze jest torusem pełnym,
a jego grupa podstawowa jest cykliczna. Nie może to
być grupa zupełnej skończenie objętościowej rozmaitości
hiperbolicznej z końcem torusowym: grupa takiego końca
jest podgrupą paraboliczną izomorficzną z $\mathbb Z^2$.
\end{uwaga}

\begin{twierdzenie}[Sztywność Mostowa--Prasada]
Niech $M_1,M_2$ będą spójnymi, zupełnymi, orientowalnymi
trójrozmaitościami hiperbolicznymi o skończonej objętości.
Jeżeli $\pi_1(M_1)\cong\pi_1(M_2)$, to istnieje izometria
$M_1\to M_2$ realizująca ten izomorfizm z dokładnością
do sprzężenia wewnętrznego. W szczególności objętość
hiperboliczna jest niezmiennikiem topologicznym.
Rozmaitości rozumiemy tutaj jako rozmaitości bez brzegu.
\end{twierdzenie}

\begin{uwaga}
Jest to zewnętrzne twierdzenie o sztywności, nie
twierdzenie wynikające jedynie z izomorfizmu grup
podstawowych dowolnych rozmaitości. Zob. G. D. Mostow,
\emph{Strong Rigidity of Locally Symmetric Spaces},
Princeton University Press, 1973; G. Prasad,
\emph{Strong rigidity of $Q$-rank $1$ lattices},
Invent. Math. 21 (1973), 255--286.
\end{uwaga}

\begin{przyklad}[Węzeł ósemkowy]
Węzeł ósemkowy $4_1$ jest hiperboliczny. Jego uzupełnienie
można przedstawić jako sklejanie dwóch regularnych idealnych
czworościanów hiperbolicznych. Stąd
$\operatorname{Vol}(4_1)=2\,v_3\approx2.0298832128$,
gdzie $v_3$ jest objętością regularnego idealnego
czworościanu w $\mathbb H^3$. Korzystamy z klasycznej
idealnej triangulacji uzupełnienia ósemki opisanej
przez Thurstona (zob. także Purcell,
\emph{Hyperbolic Knot Theory}). Ściany i krawędzie
mają miarę objętościową zero, więc objętości wnętrz
obu czworościanów dodają się. Trójlistnik nie jest
hiperboliczny, ponieważ jest węzłem torusowym.
\end{przyklad}

\begin{uwaga}[Koniec paraboliczny]
Zewnętrze hiperbolicznego węzła jest zwarte i ma brzeg
będący torusem, natomiast jego wnętrze $S^3\setminus K$
jest niezwartą rozmaitością hiperboliczną. Geometryczny
koniec ma postać $T^2\times[0,\infty)$ z metryką
$\mathrm ds^2=\mathrm dt^2+e^{-2t}g_0$,
gdzie $g_0$ jest metryką płaską na torusie.
Korzystamy z klasycznego opisu końców parabolicznych
zupełnych skończenie objętościowych rozmaitości
hiperbolicznych. W modelu półprzestrzeni taki koniec
otrzymujemy jako iloraz horokuli $z\geq z_0>0$
przez kratę translacji w płaszczyźnie $(x,y)$.
Podstawienie $z=z_0e^t$ w metryce
$z^{-2}(\mathrm dx^2+\mathrm dy^2+\mathrm dz^2)$
daje wskazaną metrykę, z
$g_0=z_0^{-2}(\mathrm dx^2+\mathrm dy^2)$ na ilorazie.
Skalowanie metryki dwuwymiarowej przez $e^{-2t}$
skaluje jej element pola przez $e^{-2t}$.
Stąd objętość końca wynosi
\[
\operatorname{Area}_{g_0}(T^2)
\int_0^\infty e^{-2t}\,\mathrm dt
=\tfrac12\operatorname{Area}_{g_0}(T^2).
\]
\end{uwaga}

\section{Chirurgia Dehna: definicja i wybór nachylenia}
\label{sec:kt-dehn}

\begin{definicja}[Nachylenie na torusie]
Niech $K\subset S^3$ będzie zorientowanym węzłem, a
$([\mu],[\lambda])$ bazą $H_1(\partial X_K;\mathbb Z)$.
\emph{Nachyleniem} na $\partial X_K$ nazywamy niezorientowaną
klasę izotopii istotnej krzywej prostej, reprezentowanej przez
pierwotny element $p[\mu]+q[\lambda]$, gdzie
$p,q\in\mathbb Z$, $\gcd(p,q)=1$, z utożsamieniem
$(p,q)\sim(-p,-q)$. Oznaczamy je symbolem $p/q$;
przez $1/0=\infty$ rozumiemy nachylenie meridianu.
\end{definicja}

\begin{definicja}[Wypełnienie i chirurgia Dehna]
Dla węzła $K\subset S^3$ i nachylenia $p/q$ wybieramy
homeomorfizm $\varphi\colon\partial(D^2\times S^1)
\to\partial X_K$ odwracający orientację brzegową
(orientacje brzegów pochodzą od zorientowanych
rozmaitości) i wysyłający brzeg dysku południkowego
$\partial D^2\times\{*\}$ na krzywą o klasie
$\pm(p[\mu]+q[\lambda])$. Istnienie takiego
sklejenia wynika z tożsamości Bézouta: pierwotny
wektor $(p,q)$ można uzupełnić do bazy całkowitej,
a przez zmianę znaku drugiego wektora uzyskać
wymagany znak wyznacznika. Macierz całkowita
o wyznaczniku $\pm1$ indukuje homeomorfizm
torusa $\mathbb R^2/\mathbb Z^2$. Rozmaitość
\[
S^3_{p/q}(K):=X_K\cup_\varphi(D^2\times S^1)
\]
nazywamy \emph{chirurgią Dehna o współczynniku $p/q$}.
Wypełnienie dla $p/q=\infty$ odtwarza $S^3$.
\end{definicja}

\begin{stwierdzenie}
Rozmaitość $S^3_{p/q}(K)$ jest dobrze określona z dokładnością
do homeomorfizmu zachowującego orientację przez nieorientowaną
klasę nachylenia $p/q$; inne wybory homeomorfizmu
$\varphi$ o zadanym obrazie meridianu torusa dołączanego
prowadzą do homeomorficznych wypełnień.
\end{stwierdzenie}

\begin{proof}
Niech $\varphi_0,\varphi_1:\partial V\to\partial X_K$
będą dwoma sklejeniami, gdzie $V=D^2\times S^1$.
Zakładamy, że obrazy klasy meridianu $m=[\partial D^2\times\{*\}]$
wyznaczają to samo niezorientowane nachylenie. Wtedy
automorfizm $h:=\varphi_0^{-1}\circ\varphi_1$ torusa
$\partial V$ zachowuje nieorientowaną klasę $m$.
Przypomnijmy elementarny opis takich homeomorfizmów.
W bazie $(m,\ell)$ macierz $h_*$ ma postać
$\begin{psmallmatrix}\varepsilon & k\\ 0 & \delta\end{psmallmatrix}$,
gdzie $\varepsilon,\delta\in\{\pm1\}$ i $k\in\mathbb Z$;
zachowanie orientacji właściwego sklejenia wymaga
$\varepsilon\delta=1$. Współrzędne brzegowe zapisujemy jako $(u,z)\in
\partial D^2\times S^1$. Dla $\varepsilon,\delta=\pm1$
niech $C_1(u)=u$ i $C_{-1}(u)=\bar u$.
Dyfeomorfizm torusa pełnego
\[
H_0(u,z)=(z^k C_\varepsilon(u),z^\delta)
\]
realizuje wskazaną macierz w bazie $(m,\ell)$.
Jest odwracalny, ponieważ z drugiej współrzędnej
odzyskujemy $z$, a następnie z pierwszej $u$.
Jego znak orientacji wynosi $\varepsilon\delta=1$.
Korzystamy z klasyfikacji klas izotopii
homeomorfizmów torusa przez $\operatorname{GL}(2,\mathbb Z)$:
$h$ jest izotopijny do ograniczenia $H_0$ do brzegu.
Izotopię różnicy przedłużamy w kołnierzu brzegu,
zatem otrzymujemy przedłużenie $H$ o dokładnie
$H|_{\partial V}=h$.
Teraz identyczność na $X_K$ i $H$ na $V$ definiują
homeomorfizm $X_K\cup_{\varphi_1}V\to
X_K\cup_{\varphi_0}V$: dla $v\in\partial V$
mamy $\varphi_0(H(v))=\varphi_1(v)$.
Oba ograniczenia zachowują orientację.
Dowód pokazuje dokładnie, dlaczego wybór całego sklejenia
nie wnosi nowej informacji poza nachyleniem meridianu.
\end{proof}


\medskip
\begin{center}
\begin{tikzpicture}[scale=.92,>=Stealth,line cap=round,line join=round]
  \fill[ktBlue!6] (0,0) rectangle (4.2,2.5);
  \draw[ktBlue,line width=1.5pt] (0,0) rectangle (4.2,2.5);
  \draw[ktInk!25,densely dashed] (1.4,0)--(1.4,2.5) (2.8,0)--(2.8,2.5)
    (0,1.25)--(4.2,1.25);
  % Zbior ilorazowy: (0,0) i (4.2,2.5) to ten sam punkt torusa.
  % Dlatego przekatna od brzegu do brzegu daje zamknieta krzywa (1,1).
  \draw[ktCoral,line width=2.8pt] (0,0)--(4.2,2.5);
  \draw[ktCoral,-{Stealth[length=3mm]},line width=2.8pt]
    (1.7,1.0119)--(2.6,1.5476);
  \node[text=ktCoral,anchor=west,font=\small] at (4.4,1.88) {$\mu+\lambda$};
  \draw[ktTeal,-{Stealth[length=2.8mm]},line width=1.6pt] (.35,-.42)--(1.55,-.42);
  \draw[ktPurple,-{Stealth[length=2.8mm]},line width=1.6pt] (-.35,.25)--(-.35,1.55);
  \node[text=ktTeal,font=\small] at (.95,-.67) {$\mu$};
  \node[text=ktPurple,font=\small] at (-.62,.92) {$\lambda$};
\end{tikzpicture}

{\footnotesize Prostokąt fundamentalny torusa $T^2=\mathbb R^2/\mathbb Z^2$. \\
Przeciwległe boki, a więc także końce czerwonej przekątnej, są utożsamione. \\
Przekątna wyznacza zamkniętą krzywą pierwotnej klasy $\mu+\lambda$ (nachylenie $1/1$).}
\end{center}
\medskip

\begin{twierdzenie}[Homologia pierwsza po chirurgii]
\label{thm:kt-dehn-h1}
Niech $K\subset S^3$ będzie węzłem, a $p,q\in\mathbb Z$
liczbami względnie pierwszymi. Wtedy
$H_1(S^3_{p/q}(K);\mathbb Z)\cong\mathbb Z/|p|\mathbb Z$
dla $p\ne0$, a dla $p=0$ jest izomorficzna z $\mathbb Z$.
Przyjmujemy $\mathbb Z/1\mathbb Z=0$.
\end{twierdzenie}

\begin{proof}
Niech $V=D^2\times S^1$ będzie dołączanym torusem pełnym,
a $\nu\in H_1(V;\mathbb Z)\cong\mathbb Z$ jego klasą rdzenia.
W ciągu Mayera--Vietorisa rozkładu
$M=X_K\cup_\varphi V$ rozważamy fragment
\[
H_1(\partial X_K)\longrightarrow H_1(X_K)\oplus H_1(V)
\longrightarrow H_1(M)\longrightarrow
H_0(\partial X_K).
\]
Następna mapa $H_0(\partial X_K)\to
H_0(X_K)\oplus H_0(V)$ jest injekcją,
gdyż wszystkie przestrzenie są spójne. Wobec tego $H_1(M)$
jest kokernem pierwszej strzałki.

Wybierzmy liczby całkowite $r,s$ spełniające
$ps-qr=1$; istnieją one, ponieważ $\gcd(p,q)=1$.
Krzywe o klasach $p\mu+q\lambda$ i
$r\mu+s\lambda$ tworzą wtedy bazę
$H_1(\partial X_K;\mathbb Z)$.
Pierwsza jest obrazem meridianu dołączonego
torusa pełnego, więc w $H_1(V)$ znika,
a w $H_1(X_K)$ ma obraz $p[\mu]$.
Druga odpowiada, po ustaleniu orientacji,
długości dołączonego torusa, zatem
jej obraz w $H_1(V)$ to $\nu$,
a w $H_1(X_K)$ to $r[\mu]$.
Dokładny fragment Mayera--Vietorisa daje więc
prezentację
\[
H_1(M;\mathbb Z)\cong
\langle[\mu],\nu\mid p[\mu]=0,\,
r[\mu]-\nu=0\rangle.
\]
Eliminujemy generator $\nu$ za pomocą
drugiej relacji i otrzymujemy
$H_1(M;\mathbb Z)\cong\mathbb Z/(p)$.
Dla $p=0$ jest to $\mathbb Z$,
a dla $p\ne0$ grupa cykliczna rzędu $|p|$.
W szczególności liczba $q$ nie wpływa na
pierwszą homologię, chociaż może zmieniać
typ homeomorficzny wypełnienia.
\end{proof}

\begin{przyklad}[Chirurgia na węźle trywialnym]
Dla węzła trywialnego $U$ jego zewnętrze $X_U$
jest torusem pełnym. Wypełnienie $p/q$ otrzymujemy
przez sklejenie dwóch torusów pełnych. Meridian
zewnętrza $X_U$ ma nachylenie $\lambda$, ponieważ
długość preferowana ogranicza w nim dysk.
Przy konwencji definiującej przestrzeń soczewkową
$L(p,q)$ właśnie przez tę operację mamy
$S^3_{p/q}(U)\cong L(p,q)$ dla $p>0$.
Dla $p<0$ można zmienić oba znaki $(p,q)$,
nie zmieniając niezorientowanego nachylenia.
Jeżeli $|p|=1$, obraz meridianu drugiego torusa
przecina meridian pierwszego algebraicznie
z modułem jeden. Skręt przedłużający się na pierwszy
torus usuwa współczynnik $q$, sprowadzając sklejenie
do meridionalnego wypełnienia $S^3$.
Dla $p=0$ oba meridiany są utożsamione.
Z jednoznaczności wypełnienia możemy wybrać
sklejenie będące iloczynem identyczności na czynniku
$S^1$ i sklejenia brzegów dwóch dysków.
Otrzymujemy $S^1\times(D^2\cup_{\partial D^2}D^2)
\cong S^1\times S^2$.
Dla $p\ne0$ twierdzenie \ref{thm:kt-dehn-h1} daje
$H_1\cong\mathbb Z/|p|\mathbb Z$.
Sama grupa $H_1$ nie rozróżnia wszystkich przestrzeni
soczewkowych; parametr $q$ może zmieniać ich topologię.
\end{przyklad}

\section{Chirurgia na splotach i macierz splecenia}
\label{sec:kt-chirurgia-link}

\begin{definicja}[Splot z ramowaniem]
Niech $L=K_1\cup\dots\cup K_r\subset S^3$ będzie
zorientowanym splotem. \emph{Ramowaniem} składowej $K_i$
nazywamy klasę homotopii nigdzie niezerowego pola normalnego
wzdłuż $K_i$, równoważnie klasę równoległej krzywej
$K_i'\subset\partial N(K_i)$. Względem długości preferowanej
$\lambda_i$ ramowanie jest całkowitą liczbą $n_i$, taką że
$[K_i']=n_i[\mu_i]+[\lambda_i]$.
Splot z ramowaniem oznaczamy $(L;n_1,\dots,n_r)$.
\end{definicja}

Ramowanie jest więc szczególnym, całkowitym nachyleniem
chirurgii. Wprowadza strukturę konieczną dla uchwytów
wymiaru cztery: bez ramowania nie można jednoznacznie
zdefiniować sposobu ich dołączenia.

\begin{definicja}[Macierz splecenia]
Dla splotu $L=K_1\cup\dots\cup K_r$ z całkowitymi
ramowaniami $n_1,\dots,n_r$ definiujemy symetryczną
macierz całkowitą $Q_L=(q_{ij})_{i,j=1}^r$ przez
$q_{ii}=n_i$ oraz $q_{ij}=\operatorname{lk}(K_i,K_j)$
dla $i\ne j$.
\end{definicja}

\begin{twierdzenie}[Homologia po integralnej chirurgii]
Niech $M_L$ będzie zamkniętą, orientowalną trójrozmaitością
otrzymaną przez integralną chirurgię na splocie
$(L;n_1,\ldots,n_r)\subset S^3$. Wtedy
\[
H_1(M_L;\mathbb Z)\cong\operatorname{coker}
(Q_L:\mathbb Z^r\longrightarrow\mathbb Z^r).
\]
W szczególności, jeśli $\det Q_L\ne0$, to
$|H_1(M_L;\mathbb Z)|=|\det Q_L|$.
\end{twierdzenie}

\begin{proof}
Niech $N_i$ będą rozłącznymi torusami pełnymi
otaczającymi składowe splotu, a $X_L$ ich wspólnym
zewnętrzem. Analogicznie jak w dowodzie stwierdzenia
\ref{prop:kt-peripheral-homology}, ciąg Mayera--Vietorisa
dla $S^3=X_L\cup\bigcup_i N_i$ daje izomorfizm
\[
\bigoplus_i H_1(\partial N_i;\mathbb Z)
\cong H_1(X_L;\mathbb Z)\oplus
\bigoplus_i H_1(N_i;\mathbb Z).
\]
Jądro projekcji indukowanej przez włączenia do $N_i$
ma bazę meridianów, zatem ich obrazy tworzą bazę
$H_1(X_L;\mathbb Z)$. Homomorfizm liczby splecenia
z $K_j$ przyjmuje na $[\mu_i]$ wartość $\delta_{ij}$.
Wyznacza więc $j$-tą współrzędną w tej bazie.
Długość preferowana $\lambda_i$ ma liczbę splecenia
zero z $K_i$, a dla $j\ne i$ tę samą liczbę splecenia
z $K_j$ co $K_i$: w $N_i$, rozłącznym z $K_j$,
obie krzywe są homologiczne. Wobec tego
\[
[\lambda_i]=\sum_{j\ne i}
\operatorname{lk}(K_i,K_j)[\mu_j]\quad\text{w }H_1(X_L;\mathbb Z).
\]

Wypełniamy wszystkie brzegi torusami pełnymi $V_i$.
W ciągu Mayera--Vietorisa dla tego rozkładu mapa
na $H_0$ jest injektywna: wektor $(a_i)_i$ przechodzi
na $(\sum_i a_i,-a_1,\ldots,-a_r)$, więc jego
zanikanie wymusza $a_i=0$ dla każdego $i$.
Homologia pierwsza wypełnienia jest zatem kokernem
mapy na $H_1$. Na $i$-tym torusie bierzemy bazę
$(n_i[\mu_i]+[\lambda_i],[\mu_i])$.
Pierwsza klasa znika w $H_1(V_i)$, druga
przechodzi na generator $\nu_i$ (dobieramy jego znak).
Otrzymujemy relacje $n_i[\mu_i]+[\lambda_i]=0$
oraz $[\mu_i]-\nu_i=0$. Eliminujemy $\nu_i$.
Pozostałe relacje są kolumnami macierzy $Q_L$
(symetria wynika z symetrii liczby splecenia),
co daje żądaną prezentację.
Jeżeli $Q_L$ jest nieosobliwa nad $\mathbb Q$, jej
kokernelem jest skończona grupa abelowa rzędu
$|\det Q_L|$, co wynika ze sprowadzenia do postaci
normalnej Smitha.
\end{proof}

\begin{uwaga}[Całkowite ramowanie i wzór na diagonalę]
Współczynnik całkowitej chirurgii na składowej $K_i$
oznacza krzywą $n_i\mu_i+\lambda_i$, gdzie $\lambda_i$
jest długością preferowaną w uzupełnieniu pojedynczego węzła
$K_i$. W zewnętrzu całego splotu jej klasa nie musi już
znikać: wynosi $\sum_{j\ne i}\operatorname{lk}(K_i,K_j)\mu_j$.
Dlatego przekątna macierzy $Q_L$ jest zadana ramowaniami,
a elementy pozadiagonalne liczbami splecenia. Nie należy
utożsamiać wartości $n_i$ z liczbą skrzyżowań w diagramie.
\end{uwaga}

\begin{przyklad}[Integralna chirurgia na splocie Hopfa]
Niech $L=K_1\cup K_2$ będzie dodatnim splotem Hopfa,
a $n_1,n_2\in\mathbb Z$ ustalonymi ramowaniami.
Wtedy $Q_L=\begin{psmallmatrix}n_1&1\\1&n_2\end{psmallmatrix}$,
$\det Q_L=n_1n_2-1$. Jeżeli $n_1=n_2=0$, to
$Q_L=\begin{psmallmatrix}0&1\\1&0\end{psmallmatrix}$,
więc $H_1(M_L;\mathbb Z)=0$. Nie oznacza to automatycznie,
że $M_L\cong S^3$ dla dowolnego splotu z taką macierzą:
homologia nie rozstrzyga typu topologicznego. Dla standardowego dodatniego splotu Hopfa można
to sprawdzić bez rachunku Kirby'ego. W modelu
$S^3\subset\mathbb C^2$ składowe to okręgi
$z_1=0$ i $z_2=0$. Ich wspólne zewnętrze
ma współrzędne $(\arg z_1,\arg z_2,|z_1|)$,
czyli jest iloczynem $T^2\times[0,1]$ po
odpowiednim przeskalowaniu ostatniej współrzędnej.
Na obu torusach brzegowych długość jednej
składowej odpowiada meridianowi drugiej.
Wypełnienie o obu współczynnikach zero
dokleja więc torusy pełne z meridianami
odpowiadającymi dwóm różnym czynnikom $T^2$.
Po usunięciu kołnierza jest to standardowe
sklejenie Heegaarda rodzaju jeden dające $S^3$,
tak samo jak w wypełnieniu o obu nachyleniach
meridionalnych (z zamienionymi rolami torusów).
\end{przyklad}

\section{Nakrycia rozgałęzione nad węzłami}
\label{sec:kt-branched}

Oprócz usuwania węzła i wykonywania chirurgii można
zbudować nową rozmaitość przez wzięcie nakrycia
rozgałęzionego. Ta operacja łączy liczbę skrzyżowań,
wyznacznik węzła i teorię form przecięcia.

\begin{definicja}[Dwukrotne nakrycie rozgałęzione]
Niech $K\subset S^3$ będzie węzłem. \emph{Dwukrotnym
nakryciem rozgałęzionym $S^3$ wzdłuż $K$} nazywamy
trójrozmaitość $\Sigma_2(K)$ z odwzorowaniem
$p\colon\Sigma_2(K)\to S^3$, które poza przeciwobrazem
$K$ jest dwukrotnym nakryciem, a w lokalnych zespolonych
współrzędnych normalnych do $K$ ma postać
$(\theta,z)\mapsto(\theta,z^2)$,
gdzie $\theta$ jest współrzędną wzdłuż węzła.
\end{definicja}

\begin{twierdzenie}[Wyznacznik z nakrycia]
Dla dowolnego węzła $K\subset S^3$ grupa
$H_1(\Sigma_2(K);\mathbb Z)$ jest skończona i
jej rząd wynosi
$|H_1(\Sigma_2(K);\mathbb Z)|
=\det(K)=|\Delta_K(-1)|$.
Jeżeli $V$ jest macierzą Seiferta węzła $K$,
to grupa ta jest izomorficzna z
$\operatorname{coker}(V+V^T:\mathbb Z^{2g}\to\mathbb Z^{2g})$.
\end{twierdzenie}

\begin{proof}[Macierz przecięć i ciąg dokładny pary]
Wybierzmy powierzchnię Seiferta $F\subset S^3$ rodzaju $g$
i przesuńmy jej wnętrze w głąb $B^4$, otrzymując właściwie
zanurzoną powierzchnię $F'\subset B^4$ o brzegu $K$.
Tworzymy dwukrotne nakrycie rozgałęzione
$W_F\to B^4$ wzdłuż $F'$. W lokalnym modelu rozgałęzienia
$z\mapsto z^2$ na normalnym dysku sprawdzamy, że brzeg
$W_F$ jest kanonicznie równy $\Sigma_2(K)$.

Powierzchnię $F$ można przedstawić jako dysk z dołączonymi
$2g$ pasmami. Opis jej nakrycia rozgałęzionego prowadzi do
modelu uchwytowego $W_F$ z jednym uchwytem indeksu zero
i $2g$ uchwytami indeksu dwa. Wynika stąd
$H_1(W_F;\mathbb Z)=0$ i $H_2(W_F;\mathbb Z)\cong\mathbb Z^{2g}$.
Przecięcia rdzeni odpowiednich uchwytów, po zamknięciu ich
powierzchniami, oblicza się przez wzajemne przesunięcia
krzywych z bazy $H_1(F)$. Diagonalne wkłady pochodzą z
ramowań pasm, a pozadiagonalne z obu kierunków przesunięcia.
Ostatecznie forma przecięcia jest reprezentowana przez
symetryzację macierzy Seiferta $S=V+V^T$.
Konstrukcję uchwytową i identyfikację formy
przecięcia przyjmujemy tutaj jako wynik zewnętrzny:
zob. S. Akbulut, R. Kirby, \emph{Branched covers of
surfaces in 4-manifolds}, Math. Ann. 252 (1980),
111--131, oraz twierdzenie Gordona--Litherlanda
o formie przecięcia nakrycia rozgałęzionego
(Invent. Math. 47 (1978), 53--69).

Ponieważ $\partial W_F=\Sigma_2(K)$, ciąg dokładny pary
i dualność Poincarégo--Lefschetza dają
\[
H_2(W_F)\xrightarrow{\;S\;}H_2(W_F,\partial W_F)
\longrightarrow H_1(\partial W_F)\longrightarrow H_1(W_F)=0.
\]
Dualność Poincarégo--Lefschetza identyfikuje
$H_2(W_F,\partial W_F;\mathbb Z)$ z $H^2(W_F;\mathbb Z)$.
Twierdzenie o uniwersalnych współczynnikach, wobec
$H_1(W_F;\mathbb Z)=0$, identyfikuje tę grupę z
$\operatorname{Hom}(H_2(W_F;\mathbb Z),\mathbb Z)$.
W bazie homologii i bazie do niej dualnej pierwsza
mapa ma więc macierz formy przecięcia $S$.
Z dokładności otrzymujemy
$H_1(\Sigma_2(K);\mathbb Z)\cong\operatorname{coker}S$.
Dla reprezentanta $P(t)=\det(tV-V^T)$
z sekcji \ref{sec:kt-alexander} mamy
$P(-1)=\det(-V-V^T)=\det S$, gdyż rozmiar
macierzy wynosi $2g$. Przy innej normalizacji
znak może się zmienić, lecz moduł pozostaje ten sam.
Ponadto $P(-1)\equiv P(1)=1\pmod 2$, zatem
wyznacznik jest niezerowy; postać normalna Smitha daje
$|H_1(\Sigma_2(K);\mathbb Z)|=|\Delta_K(-1)|$.
\end{proof}

\begin{przyklad}
Dla trójlistnika
$V+V^T=\begin{psmallmatrix}2&-1\\-1&2\end{psmallmatrix}$.
Postać normalna Smitha ma przekątną $(1,3)$,
zatem $H_1(\Sigma_2(T);\mathbb Z)\cong\mathbb Z/3\mathbb Z$.
Dla węzła ósemkowego
$V+V^T=\begin{psmallmatrix}2&-1\\-1&-2\end{psmallmatrix}$,
więc $H_1(\Sigma_2(4_1);\mathbb Z)\cong\mathbb Z/5\mathbb Z$.
Same rzędy tych grup odróżniają oba węzły.
\end{przyklad}

\begin{twierdzenie}[Homologia cyklicznych nakryć rozgałęzionych]
Niech $K\subset S^3$ będzie węzłem, a $d\ge2$ liczbą
całkowitą. Przez $\Sigma_d(K)$ oznaczmy $d$-krotne
cykliczne nakrycie $S^3$ rozgałęzione wzdłuż $K$.
Jeżeli $\zeta_d=e^{2\pi i/d}$ i
$\Delta_K(\zeta_d^j)\ne0$ dla $j=1,\dots,d-1$, to
$H_1(\Sigma_d(K);\mathbb Z)$ jest skończona i
\[
|H_1(\Sigma_d(K);\mathbb Z)|=
\left|\prod_{j=1}^{d-1}\Delta_K(\zeta_d^j)\right|.
\]
Jeżeli przynajmniej jeden z tych czynników jest zerem,
to $H_1(\Sigma_d(K);\mathbb Z)$ ma dodatnią rangę.
\end{twierdzenie}

\begin{proof}[Redukcja do modułu Alexandera]
Niech $A_K=H_1(\widetilde X_K;\mathbb Z)$ będzie modułem
Alexandera nad $\Lambda=\mathbb Z[t,t^{-1}]$.
Niech $X_d=\widetilde X_K/\langle t^d\rangle$.
Podnosimy strukturę komórkową z $X_K$ do obu nakryć.
Kompleks $C_*=C_*(\widetilde X_K;\mathbb Z)$ jest
wolny nad $\Lambda$, a kompleks $X_d$ otrzymujemy
jako $C_* /(t^d-1)C_*$. Ponieważ $t^d-1$ nie jest
dzielnikiem zera w $\Lambda$, mamy krótki ciąg
dokładny kompleksów
\[
0\longrightarrow C_*\xrightarrow{t^d-1}C_*
\longrightarrow C_*(X_d;\mathbb Z)\longrightarrow0.
\]
Nakrycie nieskończone jest spójne, więc $H_0(C_*)=\mathbb Z$,
na którym $t$ działa identycznościowo. Długi ciąg
homologii daje stąd krótki ciąg dokładny
\[
0\longrightarrow A_K/(t^d-1)A_K
\longrightarrow H_1(X_d;\mathbb Z)
\xrightarrow{\partial}\mathbb Z\longrightarrow0.
\]
Podniesienie $\mu^d$ do $X_d$ jest zamkniętą krzywą
$\widehat\mu$. Jej podniesienie do $\widetilde X_K$
ma brzeg $t^d x-x$, gdzie $x$ jest punktem startowym
(przy zgodnym wyborze generatora nakrycia).
Z definicji homomorfizmu łączącego
$\partial[\widehat\mu]=[x]=1$. Klasa ta rozszczepia
powyższy ciąg jako ciąg grup abelowych.

Brzeg $X_d$ jest jednym torusem, gdyż meridian
generuje grupę przekształceń nakrycia $\mathbb Z/d\mathbb Z$.
Jego krata podniesionych klas ma bazę
$\widehat\mu,\widehat\lambda$ odpowiadającą
$d\mu,\lambda$. Nakrycie rozgałęzione otrzymujemy
przez doklejenie torusa pełnego zgodnie z lokalnym
modelem $(\theta,z)\mapsto(\theta,z^d)$, zabijając
nachylenie $\widehat\mu$. W ciągu Mayera--Vietorisa
mapa na $H_0$ jest injektywna, a na $H_1$ druga
klasa brzegowa pozwala wyeliminować generator
rdzenia doklejonego torusa, jak w dowodzie
twierdzenia \ref{thm:kt-dehn-h1}.
W rezultacie
\[
H_1(\Sigma_d(K);\mathbb Z)
\cong H_1(X_d;\mathbb Z)/\mathbb Z[\widehat\mu]
\cong A_K/(t^d-1)A_K.
\]

Z prezentacji modułu przez macierz $tV-V^T$
z sekcji \ref{sec:kt-alexander} wynika, że ostatni iloraz jest kokernem całkowitej
macierzy otrzymanej przez podstawienie w $tV-V^T$
macierzy cyklicznego przesunięcia $d\times d$.
Po zespoleniu przesunięcie ma wartości własne
$1,\zeta_d,\dots,\zeta_d^{d-1}$. Wyznacznik
macierzy prezentacji jest więc równy
$\prod_{j=0}^{d-1}\det(\zeta_d^jV-V^T)$.
Czynnik dla $j=0$ wynosi $\det(V-V^T)=1$.
Każdy pozostały czynnik różni się od wartości
dowolnego reprezentanta $\Delta_K$ co najwyżej
o czynnik $\pm(\zeta_d^j)^k$, mający moduł jeden.
Jeżeli iloczyn jest niezerowy, postać normalna
Smitha daje rząd kokernu równy jego modułowi;
jeśli jest zerowy, macierz ma niepełny rząd nad
$\mathbb Q$, więc koker ma dodatnią rangę.

Dla $d=2$ otrzymujemy również niezależną kontrolę
prezentacji $V+V^T$. Ponieważ $\det(V-V^T)=1$,
moduł $A_K/(t-1)A_K$ jest zerowy. Wobec tego
$(t-1)A_K=A_K$ i
$(t^2-1)A_K=(t+1)A_K$. Podstawienie $t=-1$
w prezentacji modułu daje koker $-V-V^T$,
izomorficzny z kokernem $V+V^T$.
Nie korzystamy w tym rachunku z formy przecięcia
czterowymiarowego nakrycia. Otrzymujemy tezę.
\end{proof}

\begin{przyklad}
Dla trójlistnika $\Delta_T(t)=t^2-t+1$ i
$\Delta_T(-1)=3$, co daje
$|H_1(\Sigma_2(T);\mathbb Z)|=3$.
Dla $d=3$ pierwiastek $\zeta_3$ spełnia
$\zeta_3^2+\zeta_3+1=0$, więc
$\Delta_T(\zeta_3)=-2\zeta_3$ i
$\Delta_T(\zeta_3^2)=-2\zeta_3^2$.
Stąd $|H_1(\Sigma_3(T);\mathbb Z)|=4$.
\end{przyklad}

\section{Forma splecenia na wymiernej sferze homologicznej}
\label{sec:kt-linking-pairing}

\begin{definicja}[Wymierna sfera homologiczna]
Zamkniętą, spójną, orientowalną trójrozmaitość $M$
nazywamy \emph{wymierną sferą homologiczną}, jeżeli
$H_i(M;\mathbb Q)\cong H_i(S^3;\mathbb Q)$
dla wszystkich $i\geq0$. Równoważnie $H_1(M;\mathbb Z)$
jest grupą skończoną. Istotnie, homologia zwartej
trójrozmaitości jest skończenie generowana, więc
$H_1(M;\mathbb Q)=0$ zachodzi dokładnie wtedy, gdy
$H_1(M;\mathbb Z)$ jest skończona. Dualność Poincarégo
nad $\mathbb Q$ daje wtedy $H_2(M;\mathbb Q)=0$,
a spójność i orientowalność dają po jednej kopii
$\mathbb Q$ w stopniach zero i trzy.
\end{definicja}

\begin{definicja}[Forma splecenia torsyjnego]
Niech $M$ będzie orientowaną wymierną sferą homologiczną.
Dla $x,y\in H_1(M;\mathbb Z)$ wybieramy reprezentujące
je rozłączne, zorientowane 1-cykle $a,b\subset M$.
Ponieważ $x$ ma skończony rząd, istnieje $m\geq1$
i całkowity 2-łańcuch $C$ taki, że $\partial C=ma$.
Definiujemy
$\ell_M(x,y):=\frac{C\cdot b}{m}\pmod{\mathbb Z}
\in\mathbb Q/\mathbb Z$.
Tak otrzymana forma jest dobrze określona, symetryczna,
dwuliniowa i niezdegenerowana.
\end{definicja}

\begin{stwierdzenie}[Dobra określoność formy splecenia]
Dla każdej zamkniętej, spójnej, orientowanej wymiernej
sfery homologicznej $M$ liczba $\ell_M(x,y)$ z poprzedniej
definicji nie zależy od wyboru reprezentantów, liczby $m$
ani łańcucha $C$; określa niezdegenerowaną, symetryczną
formę dwuliniową na $H_1(M;\mathbb Z)$.
\end{stwierdzenie}

\begin{proof}
Niech $C,C'$ będą dwoma 2-łańcuchami z brzegiem
$m a$ (po sprowadzeniu do wspólnej wielokrotności,
co nie zmienia ilorazu). Ich różnica jest 2-cyklem.
Z dualności Poincarégo i twierdzenia o uniwersalnych
współczynnikach wynika
$H_2(M;\mathbb Z)\cong H^1(M;\mathbb Z)
\cong\operatorname{Hom}(H_1(M;\mathbb Z),\mathbb Z)=0$.
Różnica $C-C'$ jest więc brzegiem 3-łańcucha,
a jej liczba przecięcia z zamkniętym 1-cyklem
jest zero. Jeśli $a'=a+\partial E$, możemy wziąć
$C'=C+mE$; różnica ilorazów to $E\cdot b\in\mathbb Z$.
Jeżeli $b'=b+\partial D$, wzór na brzeg przecięcia
łańcuchów daje $C\cdot\partial D=(\partial C)\cdot D
=m(a\cdot D)$, więc zmiana ilorazu również jest całkowita.
Wszystkie przecięcia wykonujemy po ustawieniu
łańcuchów w położeniu ogólnym. Stąd otrzymujemy
wartość w $\mathbb Q/\mathbb Z$.

Dwuliniowość wynika z addytywności przecięć łańcuchów.
Symetrię uzyskujemy, biorąc $\partial C=ma$,
$\partial D=nb$ i porównując granicę przecięcia
$C\cap D$; daje to $n(C\cdot b)=m(a\cdot D)$,
a zatem $\ell_M(x,y)=\ell_M(y,x)$ w $\mathbb Q/\mathbb Z$.
Dla niezdegenerowania stosujemy ciąg współczynników
$0\to\mathbb Z\to\mathbb Q\to\mathbb Q/\mathbb Z\to0$.
Ponieważ $H^1(M;\mathbb Q)=H^2(M;\mathbb Q)=0$,
homomorfizm Bocksteina
\[
\beta:H^1(M;\mathbb Q/\mathbb Z)
\xrightarrow{\ \cong\ }H^2(M;\mathbb Z)
\]
jest izomorfizmem. Jeżeli $\operatorname{PD}(x)$
oznacza klasę kohomologiczną dualną do $x$, to
\[
\ell_M(x,y)=\langle\beta^{-1}\operatorname{PD}(x),y\rangle .
\]
Aby sprawdzić tę równość, reprezentujemy $x$ przez
$a$ i wybieramy $C$ z $\partial C=ma$. W kompleksie
dualnych komórek funkcjonał przecięcia z $C/m$
jest wymiernym 1-kołańcuchem. Jego kobrzeg jest
całkowitym 2-kołańcuchem przecięcia z $a$, zgodnie
ze wzorem na brzeg przecięcia użytym powyżej.
Redukcja modulo $\mathbb Z$ reprezentuje więc
$\beta^{-1}\operatorname{PD}(x)$, a obliczenie jej
na $b$ daje definicyjną wartość formy.
Uniwersalne współczynniki identyfikują
$H^1(M;\mathbb Q/\mathbb Z)$ z
$\operatorname{Hom}(H_1(M;\mathbb Z),\mathbb Q/\mathbb Z)$,
gdyż $H_0(M;\mathbb Z)=\mathbb Z$ jest wolna.
Złożenie tej identyfikacji, odwrotności $\beta$
i dualności Poincarégo jest izomorfizmem.
Jest to dokładnie odwzorowanie indukowane przez
$\ell_M$.
\end{proof}

\begin{twierdzenie}[Forma splecenia z macierzy chirurgii]
Niech $L\subset S^3$ będzie $r$-składowym splotem z
całkowitymi ramowaniami, a jego macierz splecenia
$Q\in M_r(\mathbb Z)$ niech spełnia $\det Q\ne0$.
Niech $M_L$ będzie wynikiem integralnej chirurgii na $L$.
Wówczas, przy identyfikacji
$H_1(M_L;\mathbb Z)\cong\mathbb Z^r/Q\mathbb Z^r$,
forma splecenia torsyjnego jest reprezentowana przez
$\ell_{M_L}([x],[y])\equiv -x^TQ^{-1}y\pmod{\mathbb Z}$,
przy przyjętej konwencji orientacji brzegu czterorozmaitości
uzyskanej przez dołączenie uchwytów.
\end{twierdzenie}

\begin{proof}
Dobrą określoność prawej strony sprawdzamy bezpośrednio:
po zastąpieniu $x$ przez $x+Qz$ wartość zmienia się
o $-z^Ty\in\mathbb Z$, a podobnie dla zmiany $y$.
Symetria wynika z $Q^T=Q$. Jeżeli $x^TQ^{-1}y$
jest całkowite dla każdego $y\in\mathbb Z^r$,
to $Q^{-1}x\in\mathbb Z^r$, czyli $[x]=0$.

Niech $W=W_L$ będzie śladem chirurgii, zbudowanym
z $B^4$ przez dołączenie $r$ uchwytów indeksu dwa,
i niech $M=\partial W$ z orientacją brzegową.
Kompleks uchwytowy daje $H_1(W;\mathbb Z)=0$ oraz
$H_2(W;\mathbb Z)\cong\mathbb Z^r$. W bazie
zamkniętych rdzeni forma przecięcia ma macierz $Q$:
korzystamy ze standardowego rachunku przecięć uchwytów
(zob. rozdział \ref{ch:kt-cztery}).
Przez dualność Poincarégo--Lefschetza i uniwersalne
współczynniki otrzymujemy prezentację
\[
H_2(W;\mathbb Z)\xrightarrow{Q}H_2(W,M;\mathbb Z)
\longrightarrow H_1(M;\mathbb Z)\longrightarrow0,
\]
gdzie druga baza jest dualna do pierwszej i jej
elementy reprezentują dyski współrdzeniowe, których
brzegi odpowiadają meridianom. Orientujemy je
zgodnie z tą identyfikacją.

Ustalmy $x,y\in\mathbb Z^r$. Niech $\alpha,\gamma$
będą całkowitymi 2-kocyklami na $W$ o współrzędnych
odpowiednio $x,y$ w bazie kohomologii dualnej do
zamkniętych rdzeni. Ich ograniczenia do $M$ są
dualnymi klasami kohomologicznymi do $[x],[y]$,
przez naturalność dualności względem brzegu.
Ponieważ $H^2(M;\mathbb Q)=0$, istnieje wymierny
1-kołańcuch $v$ na $M$ z $\mathrm dv=\alpha|_M$.
Jego redukcja modulo $\mathbb Z$ reprezentuje
$\beta^{-1}\operatorname{PD}([x])$. Przedłużmy $v$
dowolnie do wymiernego kołańcucha $\widetilde v$ na $W$.
Wówczas $\alpha-\mathrm d\widetilde v$ jest względnym
2-kocyklem o klasie $\widehat\alpha$ w $H^2(W,M;\mathbb Q)$,
przechodzącej na $[\alpha]$ w $H^2(W;\mathbb Q)$.
W bazie odpowiadającej $H_2(W;\mathbb Q)$ ta klasa
ma współrzędne $Q^{-1}x$. Dualność daje zatem
\[
\langle\widehat\alpha\smile[\gamma],[W,M]\rangle
=x^TQ^{-1}y.
\]

Niech $c$ będzie całkowitym 4-łańcuchem fundamentalnym
pary $(W,M)$, tak że $\partial c$ reprezentuje
klasę fundamentalną $M$. Ponieważ $\alpha,\gamma$
są całkowite, $(\alpha\smile\gamma)(c)\in\mathbb Z$.
Reguła kobrzegu iloczynu i $\mathrm d\gamma=0$ dają
\[
\begin{split}
x^TQ^{-1}y
&=(\alpha\smile\gamma)(c)
-(\mathrm d(\widetilde v\smile\gamma))(c)\\
&\equiv-(v\smile\gamma|_M)(\partial c)
\pmod{\mathbb Z}.
\end{split}
\]
Ostatnie wyrażenie jest równe
$-\ell_M([x],[y])$ na mocy powyższego opisu przez
Bocksteina i dualność. To wyprowadza również znak minus
w orientacji brzegowej (odwrócenie orientacji $M$
zmienia znak formy splecenia).
\end{proof}

\begin{przyklad}[Chirurgia na pojedynczej składowej]
Dla ramowania $p\ne0$ macierz chirurgii jest $Q=[p]$.
Stąd $H_1(M;\mathbb Z)\cong\mathbb Z/|p|\mathbb Z$,
a na generatorze $x$ zachodzi
$\ell_M(x,x)\equiv-1/p\pmod{\mathbb Z}$.
W szczególności dla przestrzeni soczewkowych sama
postać grupy homologii nie wystarcza do rozpoznania
orientowanej rozmaitości; forma splecenia przechowuje
więcej danych o topologii.
\end{przyklad}


\section{Twierdzenie Lickorisha--Wallace'a i rachunek Kirby'ego}
\label{sec:kt-wallace}

\begin{twierdzenie}[Lickorish--Wallace]
Niech $M$ będzie dowolną zamkniętą, spójną, orientowalną
trójrozmaitością. Istnieją liczba $r\geq0$, splot
$L\subset S^3$ o $r$ składowych i całkowite ramowania
$n_1,\dots,n_r$ takie, że $M$ jest homeomorficzna z zachowaniem orientacji
z trójrozmaitością otrzymaną przez integralną chirurgię
na $(L;n_1,\dots,n_r)$.
\end{twierdzenie}

\begin{proof}[Idea dowodu]
Z rozkładu Heegaarda $M=H_g\cup_\varphi H_g$ otrzymujemy
$M$ jako sklejenie dwóch brył uchwytowych wzdłuż
powierzchni rodzaju $g$. Grupę klas odwzorowań tej
powierzchni generują skręty Dehna wzdłuż odpowiednich
krzywych prostych. Porównujemy sklejenie z ustalonym rozkładem
Heegaarda sfery $S^3$ tego samego rodzaju: różnica
jest klasą odwzorowania zachowującego orientację
powierzchni. Zapisujemy ją jako skończony iloczyn
skrętów Dehna i ich odwrotności. Korzystamy z lematu
chirurgicznej realizacji skrętu: na krzywej leżącej
w powierzchni Heegaarda chirurgia o współczynniku
$+1$ albo $-1$ względem ramowania powierzchniowego
zmienia sklejenie przez odpowiedni skręt Dehna.
Umieszczamy kolejne krzywe w różnych poziomach
kołnierza powierzchni Heegaarda w $S^3$, dzięki czemu
tworzą rozłączny splot, a kolejność poziomów odpowiada
kolejności składania skrętów. Ramowanie powierzchniowe
różni się od preferowanego o liczbę całkowitą,
więc wszystkie współczynniki są całkowite.
Otrzymujemy żądane wypełnienie.
Przyjmujemy tu zewnętrznie istnienie rozkładu Heegaarda,
twierdzenie o generowaniu przez skręty i lemat
chirurgicznej realizacji skrętu. Pełny dowód wraz
z tymi składnikami: D. Rolfsen, \emph{Knots and Links},
część 9I, \emph{The Fundamental Theorem of Lickorish
and Wallace}.
\end{proof}

\begin{definicja}[Ruchy Kirby'ego dla splotów z ramowaniem]
\emph{Przesunięcie uchwytu} (handle slide) polega na zastąpieniu
jednej składowej splotu sumą pasmową z równoległą kopią
innej składowej, z odpowiednio przeniesionym ramowaniem.
\emph{Dodanie lub usunięcie wybuchu} (blow-up/down)
polega na dodaniu lub usunięciu odseparowanej od pozostałych
składowych trywialnej składowej o ramowaniu $+1$ albo $-1$,
w wersji lokalnej również z odpowiednimi zmianami sąsiednich
ramowań i skrzyżowań przy tzw. wkręceniu.
\end{definicja}

\begin{twierdzenie}[Kirby; wersja dla chirurgii integralnej]
Dwa sploty z całkowitymi ramowaniami w $S^3$ opisują
orientowanie homeomorficzne zamknięte, orientowalne
trójrozmaitości wtedy i tylko wtedy, gdy można przejść
od jednego do drugiego skończonym ciągiem izotopii,
przesunięć uchwytów oraz dodań i usunięć rozdzielnych
trywialnych składowych o ramowaniu $\pm1$.
\end{twierdzenie}

\begin{uwaga}
Jest to zewnętrzne twierdzenie R. C. Kirby'ego,
\emph{A calculus for framed links in $S^3$},
Invent. Math. 45 (1978), 35--56. Poniższy rachunek
uzasadnia niezmienniczość algebraiczną, ale nie
zupełność ruchów.
\end{uwaga}

\begin{stwierdzenie}[Algebraiczna niezmienniczość macierzy chirurgii]
\label{st:kt-kirby-matrix}
Niech $Q\in M_r(\mathbb Z)$ będzie symetryczną macierzą
splecenia splotu z całkowitymi ramowaniami.
Jeżeli przesunięcie uchwytu zastępuje ramowaną składową
$K_i$ jej sumą pasmową z równoległą kopią
$K_j$ ($i\ne j$), to nowa macierz ma postać
$Q'=P^TQP$, gdzie $P\in GL_r(\mathbb Z)$
otrzymujemy z macierzy jednostkowej przez
dodanie kolumny $j$ do kolumny $i$.
Dla sumy pasmowej z odwrotnie zorientowaną
kopią $K_j$ odejmujemy zamiast dodawać
tę kolumnę; we wszystkich przypadkach
$\det P=\pm1$.
Jeżeli dodajemy odseparowany węzeł trywialny
o ramowaniu $\varepsilon\in\{-1,1\}$,
to $Q'=Q\oplus[\varepsilon]$.

W obu przypadkach
$\operatorname{coker}Q'\cong\operatorname{coker}Q$.
Jeżeli $\det Q\ne0$, izomorfizm ten zachowuje
także postać splecenia
$\ell_Q([x],[y])=-x^TQ^{-1}y\pmod{\mathbb Z}$.
\end{stwierdzenie}

\begin{proof}
Przy przesunięciu uchwytu nowa klasa rdzenia
odpowiada algebraicznej sumie starej klasy $i$-tej
i klasy $j$-tej. Przekształcenie formy przecięcia
na uchwytach ma więc macierz $P^TQP$;
dokładniej, dla znaku $\varepsilon=\pm1$ podstawiamy
$P=I+\varepsilon E_{ji}$, gdzie $E_{ji}$ ma jedynkę
na pozycji $(j,i)$ i zera poza nią. Wówczas
$P^{-1}=I-\varepsilon E_{ji}$ i $\det P=1$.
Ramowanie nowej składowej wynosi
$q_{ii}+2\varepsilon q_{ij}+q_{jj}$,
a dla $k\ne i$ nowe splecenie wynosi
$q_{ik}+\varepsilon q_{jk}$. Pozostałe wpisy
nie zmieniają się.
Ponieważ $P$ jest unimodularna,
odwzorowanie $[x]\mapsto[P^Tx]$
z $\mathbb Z^r/Q\mathbb Z^r$
do $\mathbb Z^r/Q'\mathbb Z^r$
jest dobrze określonym izomorfizmem:
jeżeli $x=Qz$, to
$P^Tx=Q'(P^{-1}z)$.

Jeżeli $Q$ jest nieosobliwa,
mamy $(Q')^{-1}=P^{-1}Q^{-1}(P^T)^{-1}$, zatem
\[
-(P^Tx)^T(Q')^{-1}(P^Ty)
=-x^TQ^{-1}y.
\]
Dostajemy więc izometrię formy splecenia.
Dodanie składowej rozdzielnej daje macierz
$Q\oplus[\varepsilon]$, której ostatni blok
definiuje iloraz $\mathbb Z/\varepsilon\mathbb Z=0$.
W nieosobliwym przypadku ostatnia składowa
formy splecenia ma wartość całkowitą
$-\varepsilon$ i znika w $\mathbb Q/\mathbb Z$.
W obu przypadkach iloraz i forma pozostają
niezmienione.

Rachunek sprawdza \emph{niezmienniczość algebraiczną}
przy dwóch ruchach. Nie jest dowodem pełnego
twierdzenia Kirby'ego, którego część główna mówi,
że tymi ruchami można powiązać wszystkie
prezentacje tej samej orientowanej trójrozmaitości.
\end{proof}

\begin{uwaga}
Chirurgia na odseparowanej od reszty składowej trywialnej
o ramowaniu $\pm1$ nie zmienia trójrozmaitości brzegowej,
ponieważ $\pm1$-chirurgia na węźle trywialnym daje $S^3$.
Nie należy jednak mylić tego ruchu z niezmienniczością
\emph{czterorozmaitości} wypełniającej brzeg: dodanie
$\pm1$-ramowanej składowej odpowiada sumie spójnej
z $\pm\mathbb{CP}^2$ i zmienia formę przecięcia.
Właśnie to rozróżnienie będzie kluczowe w rozdziale
\ref{ch:kt-cztery}, w szczególności w sekcji \ref{sec:kt-slides}.
\end{uwaga}

\section{Hiperboliczna chirurgia Dehna}
\label{sec:kt-hyper-dehn}

\begin{twierdzenie}[Thurston; hiperboliczna chirurgia Dehna]
Niech $K\subset S^3$ będzie węzłem hiperbolicznym.
Wówczas wszystkie z wyjątkiem skończenie wielu nachyleń
$p/q$ prowadzą do zamkniętych rozmaitości
$S^3_{p/q}(K)$ dopuszczających metrykę hiperboliczną.
Ponadto dla dowolnego ciągu różnych nachyleń, których
długości w ustalonej płaskiej metryce torusa horosferycznego
dążą do nieskończoności, objętości hiperbolicznych wypełnień
dążą do $\operatorname{Vol}(K)$ od dołu
(po pominięciu skończenie wielu wyjątkowych nachyleń).
Oznacza to granicę i ścisłą nierówność, nie
koniecznie monotoniczność w zadanej kolejności.
\end{twierdzenie}

\begin{uwaga}
Przyjmujemy zewnętrzne twierdzenie o hiperbolicznym
wypełnianiu i zmianie objętości; zob. W. D. Neumann,
D. Zagier, \emph{Volumes of hyperbolic three-manifolds},
Topology 24 (1985), 307--332.
Na ustalonym torusie horosferycznym długość nachylenia
$p/q$ jest normą wektora $p u+q v$, gdzie $u,v$
tworzą bazę kraty translacji. Norma ta jest równoważna
normie euklidesowej współrzędnych $(p,q)$, więc
ograniczona długość dopuszcza tylko skończenie wiele
nachyleń. Wybierając ciąg różnych nachyleń, otrzymujemy
objętości dążące do $\operatorname{Vol}(K)$, każdą
ściśle mniejszą od tej granicy. Nie mogą przyjmować
tylko skończenie wielu wartości (wtedy ich maksimum
byłoby mniejsze od granicy). Istnieje zatem
nieskończenie wiele różnych typów homeomorficznych
wypełnień na mocy sztywności Mostowa--Prasada. Wynik stanowi pomost między geometrią
hiperboliczną a konstrukcyjną topologią trójrozmaitości.
\end{uwaga}



